\documentclass{beamer}

%\usepackage[utf8]{inputenc}

%%% MOJE
\graphicspath{{obr/}{obr/plots/}{spice/}}
\DeclareGraphicsExtensions{.ps, .eps}
%\usepackage[slovak]{babel}
\usepackage[utf8]{inputenc} 

\newcommand{\dif}{\, \mathrm{d}}	% diferencia (na derivacie)
\newcommand{\difp}{\partial}		% parc. diferencia 
\newcommand{\dxdt}[2]{\frac{\mathrm{d} #1}{\mathrm{d} #2}}
\newcommand{\dxdtp}[2]{\frac{\partial #1}{\partial #2}}
\newcommand{\un}[1]{\, \mathrm{#1}}	% jednotky velicin, v math mode
\newcommand{\E}[1]{\cdot 10^{#1}}
\newcommand{\degree}{^\circ}
\newcommand{\diameter}{\emptyset}
\newcommand{\cpx}{\widehat}		% komplexne fazory
\newcommand{\Ohm}{\Omega}

\newcommand{\myfig}[3]
{
    \begin{figure}[!ht]
    %\begin{figure}[ht]
    %\begin{figure}[htpb]
	\centering
	\includegraphics{#1}
	\caption{#2}
	%\label{fig:#3}
	#3
    \end{figure}
}
\newcommand{\myfigsc}[3]
{
    \begin{figure}[!ht]
	\centering
	%\includegraphics[width=3.5in]{#1}
	\includegraphics[scale=0.5]{#1}
	\caption{#2}
	%\label{fig:#3}
	#3
    \end{figure}
}


%\usetheme{hsrm}
%\usetheme{metropolis}
%\usecolortheme{beaver}


%Information to be included in the title page:
\title{Sample title}
\author{Anonymous}
\institute{Overleaf}
\date{2014}

\title[About Beamer] %optional
{About the Beamer class in presentation making}

\subtitle{A short story}

\author[Arthur, Doe] % (optional, for multiple authors)
{A.~B.~Arthur\inst{1} \and J.~Doe\inst{2}}

\institute[VFU] % (optional)
{
  \inst{1}%
  Faculty of Physics\\
  Very Famous University
  \and
  \inst{2}%
  Faculty of Chemistry\\
  Very Famous University
}

\date[VLC 2013] % (optional)
{Very Large Conference, April 2013}

%%%%%%%%%%%%%%%%%%%5

\title{Identification of Power BJT Operating Stages}
\subtitle{Based on Experimental Excess Charge Estimation}
\author
{J.~Miklas \and P.~Prochazka}
\institute
{\textit{Department of Power Electrical and Electronic Engineering} \\
\textit{Faculty of Electrical Engineering and Communication}\\
\textit{Brno University of Technology}\\
Brno, Czech Republic \\
}

\date{2020 IEEE 19th International Power Electronics and Motion Control Conference, Gliwice-Poland, April 2021}

\logo{\includegraphics[height=0.6cm]{logo}}
%%%%%%%%%%%%%%%%%%%5

\begin{document}

\frame{\titlepage}

\begin{frame}
\frametitle{Table of Contents}
\tableofcontents
\end{frame}


\begin{frame}
\frametitle{Sample frame title}
This is a text in the first frame. This is a text in the first frame. This is a text in the first frame.
\end{frame}
%%%%%%%%%%%%%%%%%%%%%%%%%%%%%%%%%%%%%%%%

%%%%%%%%%%%%%%%%%%%%%%%%%%%%%%%%%%%%%%%%
\begin{frame}
\frametitle{rezimy1}

\myfigsc{rezimy1_p}{Illustrative power BJT minority carrier and potential profiles for 3 basic operation modes.}{\label{fig:rezimy1}}
\end{frame}
%%%%%%%%%%%%%%%%%%%%%%%%%%%%%%%%%%%%%%%%

%%%%%%%%%%%%%%%%%%%%%%%%%%%%%%%%%%%%%%%%
\begin{frame}
\myfigsc{rezimy2}{Minority carrier concentration dependence on $I_B$ and $I_C$ variations. The geometry and concentration scales are illustrative, not proportional to real device values. Also the base and collector diagrams use the different scale to depict the same concentration gradient.}{\label{fig:rezimy2}}
\end{frame}
%%%%%%%%%%%%%%%%%%%%%%%%%%%%%%%%%%%%%%%%

%%%%%%%%%%%%%%%%%%%%%%%%%%%%%%%%%%%%%%%%
\begin{frame}
\frametitle{rezimy2}
\includegraphics[scale=0.6]{rezimy2}
%\myfigsc{rezimy1}{Illustrative power BJT minority carrier and potential profiles for 3 basic operation modes.}{\label{fig:rezimy1}}
\end{frame}
%%%%%%%%%%%%%%%%%%%%%%%%%%%%%%%%%%%%%%%%

%%%%%%%%%%%%%%%%%%%%%%%%%%%%%%%%%%%%%%%%
\begin{frame}
\frametitle{Two-column slide}

\begin{columns}

\column{0.5\textwidth}
This is a text in first column.
$$E=mc^2$$
\begin{itemize}
\item First item
\item Second item
\end{itemize}

\column{0.5\textwidth}

\myfigsc{rezimy1}{Illustrative power BJT minority carrier and potential profiles for 3 basic operation modes.}{\label{fig:rezimy1}}
\end{columns}
\end{frame}
%%%%%%%%%%%%%%%%%%%%%%%%%%%%%%%%%%%%%%%%

\end{document}
